\section{Maximum modulus and transcendence}

Fix \(R>2\).  The closed disk centered at \(2\) with radius \(R-2\) is
contained in \(\{|s|\le R\}\).  Cauchy's derivative estimate on that disk
therefore gives
\[
 |\Dpol'(2)|\le\frac{\Mdet(R)}{R-2}.
 \tag{21}
\]
Combining (20) with its certified decimal lower bound proves (7).  If
\(R\ge4\), then \(R-2\ge R/2\), which proves (8).

The trace logarithm supplies a separate qualitative conclusion.  Since
\(q_\sigma=2\alpha_0^\sigma\to0\), equation (14) gives
\(B(\sigma)\to0\).  Equations (15)--(16) therefore imply
\[
 \Dpol(\sigma)\longrightarrow1
 \qquad(\sigma\to+\infty).
 \tag{22}
\]
The derivative bound (20) proves that \(\Dpol\) is nonconstant.  A
nonconstant polynomial diverges in modulus along the positive real ray, so
(22) rules out every polynomial.  Because \(\Dpol\) is entire, it is
transcendental entire.

Finally, write \(\Dpol(s)=\sum_{m\ge0}a_ms^m\).  Transcendence means that
nonzero coefficients occur at arbitrarily high degrees.  For any fixed
\(A>0\), choose an integer \(m>A\) with \(a_m\ne0\).  Cauchy's coefficient
estimate gives
\[
 \Mdet(R)\ge |a_m|R^m,
\]
and hence
\[
 \frac{\Mdet(R)}{R^A}
 \ge |a_m|R^{m-A}\longrightarrow\infty.
\]
This proves (9), but it gives no useful numerical value for \(|a_m|\) as
\(m\) varies.
